\section{Literature review}
The simplest forms of Reinforcement Learning learn using what are known as a value functions that represents the utility of each state or state-action pair. Q-learning is one of the most commonly used of these methods. It has shown to always converge to the optimal policy as long as all states and actions are sufficiently sampled. \cite{watkins1992q}

Function approximation may be used to generalize between visited states and actions. In this case a function approximator is used to store and retrieve value function estimates. It offers a way to handle continuous state spaces, mitigate against state-action space explosion and generalize prior experience to previously unseen state-action pairs.

Deep Reinforcement Learning (DRL) is a class of RL methods where a Deep Neural Network (DNN) is used as the function approximator. As an example Deep Q-Networks (DQN) extend Q-learning by using deep learning neural networks to approximate the state-action value function. This allows agents to learn policies that can reach human-level control for high-dimensional state spaces \cite{mnih2015human}. 

An alternative to value function based methods are policy based methods. The aim is to estimate the optimal policy directly and the value is secondary if it's calculated at all. Policy gradient methods use gradient decent to estimate the policy parameters such that the expected reward is maximized. Policy based methods allow RL to be used in the continuous action spaces. An example of such a method is Deterministic Policy Gradient (DPG) \cite{silver2014deterministic}. 

Actor-critic methods are hybrid methods that combine the benefits of policy-based and value-based algorithms. In actor-critic methods the actor represents the learned policy and the critic the learned value function when its used for bootstrapping. That is the process of updating the policy and value using the estimated values of future states\cite{sutton2018reinforcement}.

Another way to use the value function in policy-based methods is to use it as a baseline when compting the parameter updates. These are known as advantage based methods. Deep Deterministic Policy Gradient (DDPG) is an extension of DPG \cite{lillicrap2015continuous} implementing both actor-critic methods and the advantage.

An intuitive way to see the advantage is that it allows agents to not only know how good the actions were, but also how much better they where than expected. Leading to reduced variance and more stable training.

One of the key issues with RL is sample efficiency. Animals are usually able to learn new tasks with just a few attempts by using prior knowledge. RL methods require too many samples before a reasonable policy is learned. This is especially an issue when collecting experience is expensive or dangerous and the problem is made worse when rewards are sparse or delayed. However several methods have been developed to mitigate these issues. 

Experience replay allows agents to remember and reuse experiences that where observed in the past, thus increases the sample efficiency. In it's simplest form experiences are replayed in the same frequency in which they ocurred regardless of the significance. An alternative is to separate the experiences into buckets of good and bad experiences and select a fixed fraction from each bucket \cite{narasimhan2015language}. Another option which has shown to be quite effective is to prioritize experiences with important transitions based on the TD error \cite{schaul2015prioritized}. 

Methods using experience replay such as DQN and DDPG have been shown to be quite effective but require large amounts of memory to store the experienced samples. In Asynchronous Advantage Actor Critic (A3C) \cite{mnih2016asynchronous} multiple instances of agents and environments are executed in parallel. Once an agent is ready it communicates back the computed gradients of the policy and value function parameters to a central parameter server and updates itself with the latest parameter values. This process is repeated periodically until the optimal policy is found. Unlike methods using experience replay which are required to use off-policy algorithms, A3C can be applied to both on and off policy algorithms.

Advantage Actor Critic (A2C) is the synchronous version of A3C. The difference is that in A2C the central parameter server waits for all agents to finish before performing the update. A2C and A3C have comparable performance \cite{mnih2016asynchronous}.

Trust Region Policy Optimization (TRPO) works by avoiding the new policy from deviating too much from the previous one, thus preventing a bad update. This is done by defining a new surrogate objective function that limits the change in policy by constraining the KL-divergence between the two policies. This method results in monotonic improvements in the policy performance, and allows for multiple updates on the policy to performed per experience sample \cite{schulman2015trust}.

Proximal Policy Optimization (PPO) works similarly to TRPO, in that it also avoids large deviation in the policy during updates. But instead of applying a hard constraint based on the KL-divergence, the surrogate objective function uses a clipping method to apply a penalty when the two policies diverge to much \cite{schulman2017proximal}.

In order to deal with sparse rewards \emph{reward shaping} allows agents to learn intermediate goals leading towards the main goal by adjusting the reward function using prior knowledge such that the agents receive rewards more frequently\cite{wiewiora2010reward}. The disadvantages of this method is that it requires prior knowledge of the environment and the optimal policy for the new adjusted reward function is not necessarily the same as the optimal policy for the original reward function. As an example Randlov \emph{et al.} \cite{randlov1998learning} where experimenting with a bicycle agent who would turn in circles to stay upright instead of reaching its goal.

Difference (D) \cite{wolpert2000collective} rewards and potential-based reward shaping (PBRS) \cite{ng1999policy}, are commonly used shaping approaches which have been successfully applied to a wide range of domains. 

Another method for handling sparse rewards is Learning from Demonstrations (LfD). In LfD, an agent learns to perform a new task from demonstrations performed by an expert without feedback from the rewards. Obtaining sufficient high quality demonstrations may not be feasible in many cases, thus leading to suboptimal policies. However LfD is still a useful tool that can be used to initialize the agent with a good or a safe policy and further learning can then search for better policies. 

For example AlphaGo \cite{silver2016mastering} initializes the policy network using supervised learning on state-action pairs from recorded human experts. Then using self-play AlphaGo was able to discover new plays and learn from its mistakes to eventually achieve super human performance.

Another approach is that taken by Hester \emph{et al.} \cite{hester2018deep} to develop Deep Q-learning from Demonstrations (DQfD). In this case the agent is initialized by placing a small set of demonstrations in the replay buffer. DQfD pre-trains solely on the demonstration data using a combination of both TD and supervised losses and eventually the agent starts learning using a combination of both demonstrated and self-generated data. This resulted in quite a capable method that outperformed standard reinforcement learning and three related methods for incorporating demonstration data into DQN.

Behavior Cloning (BC) is a class supervised learning methods that maps states to actions based on demonstrations from an expert. On the other hand Inverse Reinforcement Learning (IRL) estimates the underlying value function that justifies the actions taken by the expert. This is considered to be more robust.

A drawback of IRL algorithms is that they can be computationally expensive since they require a reinforcement learning inner loop between the cost estimation, policy training and evaluation. 

An alternative to IRL inspired by Generative Adversarial Nets (GAN) is Generative
Adversarial Imitation Learning. It avoids the use of the inner loop by training a policy until it is able to fool a discriminator, resulting in policies that travel the same MDP paths as the expert.
